% !TEX root = main.tex
\chapter{Automorphism Groups}
\label{chap:automorphism-groups}

Fraïssé theory leads naturally from countable structures to their automorphism groups. In this chapter we study automorphism groups of first-order structures as permutation groups on the underlying domain and equip them with the topology of pointwise convergence. We begin with the basic notions of stabilizers, orbits, transitivity, and rigidity, and then introduce the permutation topology on $\Sym(\Omega)$ and on its subgroups.

We next characterize those subgroups of $\Sym(\Omega)$ which arise as automorphism groups of structures on $\Omega$, in terms of closed subgroups and canonical structures. In the countable case, this shows that automorphism groups are closed subgroups of $S_\infty$, and hence naturally carry the structure of Polish groups. These results provide the topological framework needed in the following chapters, where we study continuous actions of such groups on compact spaces.

\section{Automorphism Groups as Permutation Groups}

Let $A$ be an $L$-structure with domain $\Omega=\dom(A)$. Every automorphism of $A$ is a permutation of $\Omega$, and the set of all automorphisms of $A$ forms a group under composition. We denote this group by $\Aut(A)$ and regard it as a subgroup of the symmetric group $\Sym(\Omega)$.

More generally, if $\Omega$ is any set, we write $\Sym(\Omega)$ for the group of all permutations of $\Omega$. Thus every automorphism group is naturally a permutation group. Our aim in this chapter is to study those subgroups of $\Sym(\Omega)$ which arise as automorphism groups of structures on $\Omega$, and to equip them with the topology of pointwise convergence.

Let $G\leq\Sym(\Omega)$. For a subset $X\subseteq\Omega$, the \emph{pointwise stabilizer} of $X$ in $G$ is \(G_{(X)}=\{g\in G:g(a)=a\text{ for all }a\in X\}\), and the \emph{setwise stabilizer} of $X$ in $G$ is \(G_{\{X\}}=\{g\in G:g(X)=X\}\).
Thus $G_{(X)}\leq G_{\{X\}}\leq G$. If $\bar a=(a_1,\dots,a_n)$ is a finite tuple from $\Omega$, we also write $G_{(\bar a)}$ for the pointwise stabilizer of the set $\{a_1,\dots,a_n\}$.

For $a\in\Omega$, the \emph{orbit} of $a$ under $G$ is the set \(G\cdot a=\{g(a):g\in G\}\). More generally, if $\bar a=(a_1,\dots,a_n)\in\Omega^n$, the \emph{orbit} of $\bar a$ under $G$ is \(G\cdot \bar a=\{g\bar a:g\in G\}\), where \(g\bar a=(g(a_1),\dots,g(a_n))\).
The orbits of $G$ on $\Omega$ partition $\Omega$. We say that $G$ is \emph{transitive} if it has exactly one orbit on $\Omega$.

Accordingly, a structure $A$ is called \emph{transitive} if $\Aut(A)$ acts transitively on $\dom(A)$. At the opposite extreme, $A$ is called \emph{rigid} if $\Aut(A)=\{1_{\Omega}\}$.

\begin{example}
    The following examples will be useful throughout the thesis.
    \begin{enumerate}[(i)]
        \item If $A$ is a pure set with no non-logical symbols, then every permutation of $\dom(A)$ is an automorphism, so $\Aut(A)=\Sym(\dom(A))$.
        \item Every finite linear order is rigid. Indeed, if $A=(\{a_1,\dots,a_n\},<)$ with $a_1<\cdots<a_n$, then any automorphism must preserve the least element, then the next least element, and so on; hence it fixes every point.
        \item The ordered set $(\mathbb Q,<)$ is transitive. Given $q,r\in\mathbb Q$, there exists an order automorphism of $(\mathbb Q,<)$ sending $q$ to $r$, so all elements lie in a single orbit under $\Aut(\mathbb Q,<)$.
    \end{enumerate}
\end{example}

These notions will reappear throughout the chapter. In particular, stabilizers of finite tuples will provide the basic open neighbourhoods in the permutation topology, while orbits on finite powers $\Omega^n$ will play a central role in the description of closed and dense subgroups.

\section{The Permutation Topology}

We now introduce the natural topology on permutation groups. Let $\Omega$ be a set. For $n<\omega$ and tuples $\bar a,\bar b\in\Omega^n$, set \(S(\bar a,\bar b)=\{g\in\Sym(\Omega):g\bar a=\bar b\}\).
A subset of $\Sym(\Omega)$ is called \emph{basic open} if it is of the form $S(\bar a,\bar b)$ for some finite tuples $\bar a,\bar b$ of the same length. An \emph{open} subset of $\Sym(\Omega)$ is any union of basic open sets. This is the \emph{permutation topology}, or equivalently the topology of pointwise convergence.

If $G\leq\Sym(\Omega)$, we always equip $G$ with the subspace topology induced from $\Sym(\Omega)$. Thus the basic open subsets of $G$ are precisely the sets \(G\cap S(\bar a,\bar b)\), where $\bar a,\bar b$ are finite tuples from $\Omega$ of the same length.

\begin{proposition}
    \label{prop:permutation-topology}
    Let $G\leq\Sym(\Omega)$.
    \begin{enumerate}[(a)]
        \item The sets $S(\bar a,\bar b)$ form a basis for a topology on $\Sym(\Omega)$. With this topology, $\Sym(\Omega)$ is a topological group, and hence so is every subgroup $G\leq\Sym(\Omega)$ with the induced topology.
        \item A subgroup $H\leq G$ is open if and only if it contains $G_{(\bar a)}$ for some finite tuple $\bar a$ from $\Omega$.
        \item A subset $F\subseteq G$ is closed if and only if whenever $g\in G$ has the property that every basic open neighbourhood of $g$ meets $F$, then $g\in F$.
        \item A subgroup $H\leq G$ is dense in $G$ if and only if $H$ and $G$ have the same orbits on $\Omega^n$ for every positive integer $n$.
    \end{enumerate}
\end{proposition}

\begin{proof}
    For (a), note first that \(\Sym(\Omega)=S(\emptyset,\emptyset)\).
    Moreover, if $g\in S(\bar a,\bar b)\cap S(\bar c,\bar d)$, then
    \[
        g\in S(\bar a\bar c,\bar b\bar d)\subseteq S(\bar a,\bar b)\cap S(\bar c,\bar d),
    \]
    where $\bar a\bar c$ and $\bar b\bar d$ denote concatenation of tuples. Thus the sets $S(\bar a,\bar b)$ form a basis for a topology on $\Sym(\Omega)$.

    To prove continuity of inversion, observe that \(g\in S(\bar a,\bar b)\Longleftrightarrow g^{-1}\in S(\bar b,\bar a)\). For continuity of multiplication, let $(g,h)\in\Sym(\Omega)\times\Sym(\Omega)$ and suppose that \(gh\in S(\bar a,\bar b)\).
    Set $\bar c=h\bar a$. Then
    \[
        h\in S(\bar a,\bar c)\qquad\text{and}\qquad g\in S(\bar c,\bar b),
    \]
    and whenever $g'\in S(\bar c,\bar b)$ and $h'\in S(\bar a,\bar c)$, we have \(g'h'\bar a=g'\bar c=\bar b\), so that \(g'h'\in S(\bar a,\bar b)\).
    Hence
    \[
        S(\bar c,\bar b)\cdot S(\bar a,\bar c)\subseteq S(\bar a,\bar b),
    \]
    which proves continuity of multiplication. Therefore $\Sym(\Omega)$ is a topological group. Since subgroups inherit the subspace topology, every subgroup $G\leq\Sym(\Omega)$ is also a topological group.

    For (b), if $\bar a$ is a finite tuple from $\Omega$, then \(G_{(\bar a)}=G\cap S(\bar a,\bar a)\),
    so $G_{(\bar a)}$ is open in $G$. If $H\leq G$ contains $G_{(\bar a)}$, then
    \[
        H=\bigcup_{h\in H} hG_{(\bar a)},
    \]
    and each left coset $hG_{(\bar a)}$ is open, since left translation is a homeomorphism of $G$. Thus $H$ is open.

    Conversely, suppose that $H\leq G$ is open. Since $1_G\in H$, there exists a basic open neighbourhood of $1_G$ contained in $H$. Such a neighbourhood has the form \(G\cap S(\bar a,\bar a)=G_{(\bar a)}\) for some finite tuple $\bar a$ from $\Omega$. Hence $G_{(\bar a)}\subseteq H$.

    Statement (c) is just the usual characterization of closed sets in terms of neighbourhoods: a subset $F\subseteq G$ is closed if and only if every point outside $F$ has an open neighbourhood disjoint from $F$, equivalently, whenever every basic open neighbourhood of $g$ meets $F$, one has $g\in F$.

    For (d), first suppose that $H$ is dense in $G$, and let $\bar a,\bar b\in\Omega^n$ lie in the same $G$-orbit. Choose $g\in G$ with $g\bar a=\bar b$. Then $S(\bar a,\bar b)$ is a basic open neighbourhood of $g$, so by density it meets $H$. Hence there exists $h\in H$ such that $h\bar a=\bar b$. Thus $\bar a$ and $\bar b$ lie in the same $H$-orbit. It follows that every $G$-orbit on $\Omega^n$ is also an $H$-orbit, and therefore $H$ and $G$ have the same orbits on $\Omega^n$.

    Conversely, suppose that $H$ and $G$ have the same orbits on $\Omega^n$ for every $n\geq 1$. To show that $H$ is dense in $G$, it suffices by part (c) to show that every nonempty basic open subset of $G$ meets $H$. Let \(G\cap S(\bar a,\bar b)\) be a nonempty basic open subset of $G$. Then there exists $g\in G$ such that $g\bar a=\bar b$, so $\bar a$ and $\bar b$ lie in the same $G$-orbit on $\Omega^n$. By assumption they also lie in the same $H$-orbit, so there exists $h\in H$ such that $h\bar a=\bar b$. Hence \(h\in H\cap S(\bar a,\bar b)\subseteq H\cap G\cap S(\bar a,\bar b)\), and therefore $G\cap S(\bar a,\bar b)$ meets $H$. This proves that $H$ is dense in $G$.
\end{proof}

In particular, the stabilizers of finite tuples form a neighbourhood basis at the identity in any subgroup of $\Sym(\Omega)$, and orbit structure on finite powers $\Omega^n$ detects density. These two observations will be used repeatedly in the next section.

\section{Closed Subgroups and Canonical Structures}

We now characterize those subgroups of $\Sym(\Omega)$ which arise as automorphism groups of structures on $\Omega$. In view of the preceding section, the relevant notion of closedness is topological closedness in the permutation topology.

\begin{theorem}
    \label{thm:closed-subgroups-automorphism-groups}
    Let $\Omega$ be a set, let $G\leq \Sym(\Omega)$, and let $H\leq G$. Then the following are equivalent.
    \begin{enumerate}[(a)]
        \item $H$ is closed in $G$.
        \item There exists a structure $A$ with domain $\Omega$ such that
              \[
                  H=G\cap \Aut(A).
              \]
    \end{enumerate}
    In particular, a subgroup $H\leq \Sym(\Omega)$ is of the form $\Aut(A)$ for some structure $A$ with domain $\Omega$ if and only if $H$ is closed in $\Sym(\Omega)$.
\end{theorem}

\begin{proof}
    Assume first that $H$ is closed in $G$. For each $n<\omega$ and each orbit $\Delta$ of $H$ on $\Omega^n$, introduce an $n$-ary relation symbol $R_\Delta$. Let $L$ be the language consisting of all these relation symbols, and define an $L$-structure $A$ with domain $\Omega$ by setting \(R_\Delta^A=\Delta\) for every orbit $\Delta$ of $H$ on $\Omega^n$.

    Every element of $H$ preserves each relation $R_\Delta^A$, since each $R_\Delta^A$ is an $H$-orbit. Hence \(H\subseteq G\cap \Aut(A)\).

    Conversely, let $g\in G\cap \Aut(A)$. We claim that every basic open neighbourhood of $g$ in $G$ meets $H$. Let $\bar a\in\Omega^n$, and let $\Delta$ be the $H$-orbit of $\bar a$. Since $\bar a\in \Delta=R_\Delta^A$ and $g$ is an automorphism of $A$, we have
    \[
        g\bar a\in R_\Delta^A=\Delta.
    \]
    Thus there exists $h\in H$ such that \(h\bar a=g\bar a\). It follows that \(h\in H\cap G\cap S(\bar a,g\bar a)\),
    so the basic open neighbourhood $G\cap S(\bar a,g\bar a)$ of $g$ meets $H$. Since $H$ is closed in $G$, we conclude that $g\in H$. Therefore \(G\cap \Aut(A)\subseteq H\), and hence \(H=G\cap \Aut(A)\).

    Now assume that there exists a structure $A$ with domain $\Omega$ such that \(H=G\cap \Aut(A)\).
    To show that $H$ is closed in $G$, let $g\in G$ and suppose that every basic open neighbourhood of $g$ in $G$ meets $H$. We show that $g\in \Aut(A)$.

    Let $\varphi(\bar x)$ be an atomic formula in the language of $A$, and let $\bar a$ be a tuple from $\Omega$ of the same length as $\bar x$. By assumption, the basic open neighbourhood \(G\cap S(\bar a,g\bar a)\) meets $H$. Choose \(h\in H\cap G\cap S(\bar a,g\bar a)\).
    Then $h\bar a=g\bar a$. Since $h\in H\subseteq \Aut(A)$, we have
    \[
        A\models \varphi(\bar a)\iff A\models \varphi(h\bar a).
    \]
    Since $h\bar a=g\bar a$, it follows that \(A\models \varphi(\bar a)\iff A\models \varphi(g\bar a)\). Thus $g$ preserves all atomic formulas, and since $g$ is bijective it follows that \(g\in G\cap \Aut(A)=H\).
    Hence $H$ is closed in $G$.
\end{proof}

\begin{corollary}
    \label{cor:automorphism-groups-are-closed}
    Let $A$ be a structure with domain $\Omega$. Then $\Aut(A)$ is a closed subgroup of $\Sym(\Omega)$.
\end{corollary}

\begin{proof}
    Apply Theorem~\ref{thm:closed-subgroups-automorphism-groups} with $G=\Sym(\Omega)$ and $H=\Aut(A)$.
\end{proof}

When $H$ is a closed subgroup of $\Sym(\Omega)$, the structure constructed in the proof of Theorem~\ref{thm:closed-subgroups-automorphism-groups} is called the \emph{canonical structure} of $H$. Its domain is $\Omega$, and for each $n<\omega$ its basic $n$-ary relations are precisely the $H$-orbits on $\Omega^n$. By construction, \(H=\Aut(A)\).

Thus the canonical structure records exactly the orbit structure of $H$ on finite powers of $\Omega$. In particular, closed subgroups of $\Sym(\Omega)$ and automorphism groups of structures on $\Omega$ may be regarded as two equivalent ways of encoding the same data.

\section{Countable Automorphism Groups as Polish Groups}

We now specialize to the countable case, which is the one relevant for Fra\"{\i}ss\'e limits. Throughout this section, let $\Omega$ be a countable set, and fix an enumeration \(\Omega=\{\omega_1,\omega_2,\dots\}\).
Our goal is to show that automorphism groups of countable structures are Polish. Since $\Sym(\Omega)$ is already a topological group by Proposition~\ref{prop:permutation-topology}, it remains to verify in this setting separability, complete metrizability, and the fact that closed subgroups of Polish groups are again Polish.


\begin{proposition}
    \label{prop:sym-countable-second-countable}
    The permutation topology on $\Sym(\Omega)$ is second countable. Consequently, every subgroup of $\Sym(\Omega)$, endowed with the induced topology, is also second countable.
\end{proposition}

\begin{proof}
    For each $n<\omega$ and each pair of tuples $\bar a,\bar b\in\Omega^n$, the set $S(\bar a,\bar b)$ is basic open. Since $\Omega$ is countable, each $\Omega^n$ is countable, and hence
    \[
        \bigcup_{n<\omega}\Omega^n\times\Omega^n
    \]
    is countable. Therefore there are only countably many basic open sets of the form $S(\bar a,\bar b)$, so $\Sym(\Omega)$ has a countable basis. Since second countability passes to subspaces, every subgroup of $\Sym(\Omega)$ is also second countable.
\end{proof}

\begin{corollary}
    \label{cor:sym-countable-separable}
    The space $\Sym(\Omega)$ is separable. Consequently, every subgroup of $\Sym(\Omega)$ is separable.
\end{corollary}

\begin{proof}
    Every second countable space is separable, so the conclusion follows from Proposition~\ref{prop:sym-countable-second-countable}.
\end{proof}

We now describe an explicit compatible complete metric. First define a metric $d$ on $\Sym(\Omega)$ by
\[
    d(g,h)=
    \begin{cases}
        0,            & \text{if } g=h,                                                                   \\[4pt]
        \dfrac{1}{n}, & \text{if } g\neq h,\ \text{where } n=\min\{m\geq 1:g(\omega_m)\neq h(\omega_m)\}.
    \end{cases}
\]
Thus $d(g,h)$ records the first point of the enumeration on which $g$ and $h$ disagree.

\begin{proposition}
    \label{prop:d-compatible}
    The function $d$ is a metric on $\Sym(\Omega)$, and it induces the permutation topology.
\end{proposition}

\begin{proof}
    The only nontrivial point is the triangle inequality. Let $g,h,k\in\Sym(\Omega)$, and suppose that
    \[
        d(g,k)=\frac{1}{a}, \qquad d(k,h)=\frac{1}{b}.
    \]
    Then $g(\omega_m)=k(\omega_m)=h(\omega_m)$ for every $m<\min(a,b)$. Hence
    \[
        d(g,h)\leq \frac{1}{\min(a,b)}=\max\{d(g,k),d(k,h)\},
    \]
    so $d$ is in fact an ultrametric.

    For $g\in\Sym(\Omega)$ and $n\geq 1$, the open ball of radius $1/(n+1)$ about $g$ is
    \[
        B_d\!\left(g,\frac{1}{n+1}\right)
        =
        \{h\in\Sym(\Omega):h(\omega_m)=g(\omega_m)\text{ for all }m\leq n\}.
    \]
    Thus the $d$-balls are exactly the sets of permutations agreeing with $g$ on a finite initial segment of the enumeration. These form a neighbourhood basis at $g$ for the permutation topology, so $d$ induces that topology.
\end{proof}

The metric $d$ is not complete, since a sequence of longer and longer finite cycles can be Cauchy for $d$ while converging pointwise to an injective non-surjective map. To obtain a complete compatible metric, define \(d'(g,h)=d(g,h)+d(g^{-1},h^{-1})\).

\begin{proposition}
    \label{prop:dprime-compatible-complete}
    The function $d'$ is a complete metric on $\Sym(\Omega)$, and it induces the permutation topology.
\end{proposition}

\begin{proof}
    Since $d$ is a metric, so is $d'$. Moreover, \(d(g,h)\leq d'(g,h)\),
    so every $d'$-open set is $d$-open. Conversely, if $h$ agrees with $g$ on $\omega_1,\dots,\omega_n$, then every permutation sufficiently close to $h$ in the metric $d'$ also agrees with $g$ on those points. Hence every basic $d$-open neighbourhood is $d'$-open, and the two metrics induce the same topology.

    It remains to prove completeness. Let $F=\Omega^\Omega$, and define a metric $\delta$ on $F$ by
    \[
        \delta(f_1,f_2)=
        \begin{cases}
            0,            & \text{if } f_1=f_2,                                                                       \\[4pt]
            \dfrac{1}{n}, & \text{if } f_1\neq f_2,\ \text{where } n=\min\{m\geq 1:f_1(\omega_m)\neq f_2(\omega_m)\}.
        \end{cases}
    \]
    Exactly as above, $(F,\delta)$ is complete. Equip $F\times F$ with the product metric
    \[
        \rho\bigl((f_1,h_1),(f_2,h_2)\bigr)=\delta(f_1,f_2)+\delta(h_1,h_2).
    \]
    Then $(F\times F,\rho)$ is complete.

    Now define \(\Phi:\Sym(\Omega)\to F\times F\) by \(\Phi(g)=(g,g^{-1})\). For $g,h\in\Sym(\Omega)$,
    \[
        \rho(\Phi(g),\Phi(h))
        =\delta(g,h)+\delta(g^{-1},h^{-1})
        =d'(g,h),
    \]
    so $\Phi$ is an isometric embedding.

    Its image consists exactly of the pairs $(f,h)\in F\times F$ satisfying \(f\circ h=\operatorname{id}_{\Omega}\) and \(h\circ f=\operatorname{id}_{\Omega}\). This is a closed subset of $F\times F$, since for each $\omega\in\Omega$ the conditions \(f(h(\omega))=\omega\) and \(h(f(\omega))=\omega\) define closed sets. Therefore $\Phi(\Sym(\Omega))$ is complete, and hence $(\Sym(\Omega),d')$ is complete.
\end{proof}

\begin{definition}
    A \emph{Polish group} is a topological group whose topology is separable and completely metrizable.
\end{definition}

\begin{proposition}
    \label{prop:closed-subgroup-polish}
    Let $G$ be a Polish group and let $H\leq G$ be a closed subgroup. Then $H$, equipped with the subspace topology, is again a Polish group.
\end{proposition}

\begin{proof}
    Since multiplication and inversion on $H$ are restrictions of the corresponding continuous maps on $G$, the group $H$ is a topological group with the subspace topology.

    Let $d$ be a compatible complete metric on $G$. Then the restriction $d|_{H\times H}$ is a compatible metric on $H$, and it is complete because $H$ is closed in the complete metric space $(G,d)$.

    Since $G$ is Polish, it is in particular separable and metrizable, hence second countable. Therefore the subspace $H$ is also second countable, and thus separable. It follows that $H$ is separable and completely metrizable, so $H$ is Polish.
\end{proof}
\begin{theorem}
    \label{thm:aut-countable-polish}
    Let $A$ be a structure with countable domain $\Omega$. Then $\Aut(A)$, endowed with the permutation topology, is a Polish group.
\end{theorem}

\begin{proof}
    By Corollary~\ref{cor:automorphism-groups-are-closed}, the group $\Aut(A)$ is a closed subgroup of $\Sym(\Omega)$. By Proposition~\ref{prop:dprime-compatible-complete}, the group $\Sym(\Omega)$ is completely metrizable, and by Corollary~\ref{cor:sym-countable-separable} it is separable. Since $\Sym(\Omega)$ is a topological group by Proposition~\ref{prop:permutation-topology}, it is therefore a Polish group. Now Proposition~\ref{prop:closed-subgroup-polish} shows that $\Aut(A)$ is Polish.
\end{proof}

\begin{proposition}
    \label{prop:countable-local-basis-open-subgroups}
    Let $G\leq\Sym(\Omega)$. Then the family
    \[
        \{G_{(F)}:F\subseteq \Omega \text{ finite}\}
    \]
    is a countable neighbourhood basis at the identity consisting of open subgroups.
\end{proposition}

\begin{proof}
    Each $G_{(F)}$ is a subgroup of $G$, and by Proposition~\ref{prop:d-compatible} it is open. By Proposition~\ref{prop:sym-countable-second-countable}, the finite subsets of $\Omega$ form a countable family. Finally, every basic open neighbourhood of the identity in $G$ contains $G_{(F)}$ for some finite set $F\subseteq\Omega$, so these subgroups form a neighbourhood basis at the identity.
\end{proof}

In particular, automorphism groups of countable structures are non-archimedean Polish groups: they admit a countable local basis at the identity consisting of open subgroups. This is the topological setting in which we will study flows, minimality, and universal minimal flows in the following chapters.
